% ─────────────────────────────────────────────────────────────────────────────
\section{Finite Simple Continued Fractions}
% ─────────────────────────────────────────────────────────────────────────────

We now turn to the case of finite simple continued fractions, which correspond to rational numbers.
As mentioned in the definition of the \texttt{SimpleContinuedFraction} class, the body generator 
$g$ is allowed to have a finite domain, in which case the resulting SCF is finite.

Every rational $\nicefrac{p}{q}$ number admits exactly two finite SCF representations, one of which has a final partial quotient of 1.
\[
    \frac{p}{q} = [a_0; a_1, \ldots, a_n] = [a_0; a_1, \ldots, a_{n-1}, a_n - 1, 1] \quad \text{with } a_n \geq 2.
\]

When it is said that a rational number has a \textit{unique} SCF representation, it is understood
that the representation is taken in canonical form, i.e. with $a_n \geq 2$. While this condition
is not enforced at runtime, all constructors that produce finite SCFs (e.g. 
\texttt{FiniteContinuedFraction.from\_rational}) return the canonical form.

\subsection{The Euclidean Algorithm as SCF Expansion}

The SCF coefficients of $p/q$ are precisely the successive quotients produced by the
Euclidean algorithm.
Concretely, \texttt{Convert.from\_rational\_to\_scf} implements:
\begin{align*}
  p_0 &= p, \quad p_1 = q, \\
  a_k &= \left\lfloor \frac{p_k}{p_{k+1}} \right\rfloor, \\
  p_{k+2} &= p_k \bmod p_{k+1}, \qquad k \geq 0,
\end{align*}
until a zero remainder is reached.

\begin{remark}
This correspondence between the Euclidean algorithm and SCF expansion provides
the bijection between finite (canonical) SCFs and rational numbers.

The library uses \texttt{divmod} (Python's built-in, backed by hardware division) for each step,
relying on Python's arbitrary-precision integers so that very large rationals are handled correctly
without overflow.
\end{remark}

\begin{catenadef}[Conversion from Rational-like Types to Finite SCF]
\begin{lstlisting}[style=python]
from catena import FiniteSimpleContinuedFraction

# Convert a rational (Fraction | tuple[int, int] | int) number to its finite SCF representation.
rational_number = (17, 15) # Represents 17/15
finite_scf = FiniteSimpleContinuedFraction.from_rational(rational_number)
print(finite_scf) # FiniteSimpleContinuedFraction(partial_quotients=(7, 2), integer_part=1)

# Convert a float number to its finite SCF representation.
# Note: this conversion may not be exact due to the nature of floating-point representation.
float_number = 3.14
finite_scf_from_float = FiniteSimpleContinuedFraction.from_float(float_number)
print(finite_scf_from_float) # FiniteSimpleContinuedFraction(partial_quotients=(7, 7, 3216856876693, 14), integer_part=3)

finite_scf_from_float_limited = FiniteSimpleContinuedFraction.from_float(3.14, max_denominator=10000) # Limit the denominator to 10,000 to elude floating-point precision artifacts.
print(finite_scf_from_float_limited) # FiniteSimpleContinuedFraction(partial_quotients=(7, 7), integer_part=3)

# Convert a decimal string to its finite SCF representation.
decimal_string = "3.14"
finite_scf_from_decimal = FiniteSimpleContinuedFraction.from_decimal(decimal_string)
print(finite_scf_from_decimal) # FiniteSimpleContinuedFraction(partial_quotients=(7, 7), integer_part=3)
\end{lstlisting}
\end{catenadef}

\begin{remark}
  \texttt{FiniteSimpleContinuedFraction.from\_decimal} takes a decimal string as input rather than a 
  \texttt{decimal.Decimal} object. This is because the latter depends on a global precision context 
  (\texttt{getcontext().prec}), and operations such as \texttt{Decimal(2).sqrt()} produce approximations 
  whose accuracy depends on that context. As a result, converting directly from \texttt{Decimal} would 
  require an additional precision parameter and would not be guaranteed to produce consistent results 
  across different environments. Using a decimal string avoids this ambiguity by making the input 
  representation explicit in the API call.
\end{remark}